\documentclass[11pt]{article}
\usepackage{mathblatt}


\begin{document}
\pruefheftstil{Wachstum}{Prüfungsheft P10}
\noindent{\Large\bfseries Wachstum}\hfill{\small\color{mbgrau}Prüfungsheft P10}\par\medskip
\pfrueckblick
\begin{pfnr}{}{1.}{}
Ein Ticket kostet $50\,€$. Der Preis steigt um $10\,\%$. Wie hoch ist der neue Preis?
\par\noindent \leerfeld[€] \par
\rechenplatz{2}
\fusshilfe{\mbox{1:~$55\,€$}}
\end{pfnr}
\begin{pfnr}{}{2.}{}
Die Tabelle zeigt Werte von $x$ und $y$. Welcher $y$-Wert gehört zu $x = 2$? \wertetabelle[5,8,11,14]{x}{y}{0,1,2,3}
\par\noindent y = \leerfeld \par
\rechenplatz{2}
\fusshilfe{\mbox{2:~$11$}}
\end{pfnr}
\begin{pfnr}{}{3.}{}
Berechne $1{,}5 \cdot 4$.
\par\noindent \leerfeld \par
\rechenplatz{1}
\fusshilfe{\mbox{3:~$6$}}
\end{pfnr}
\begin{pfnr}{}{4.}{}
Wie viel sind $20\,\%$ von $150\,€$?
\par\noindent \leerfeld[€] \par
\fusshilfe{\mbox{4:~$30\,€$}}
\end{pfnr}
\begin{pfnr}{}{5.}{}
In einer Tabelle steht der Wert $45$. Darunter steht der Wert $63$. Berechne den Quotienten $\tfrac{63}{45}$.
\par\noindent \leerfeld \par
\rechenplatz{1}
\fusshilfe{\mbox{5:~$1{,}4$}}
\end{pfnr}
\begin{pfnr}{}{6.}{}
Berechne $2^5$ mit dem Taschenrechner.
\par\noindent \leerfeld \par
\rechenplatz{1}
\fusshilfe{\mbox{6:~$32$}}
\end{pfnr}
\pfstufekopf[stabelleergnzen]{Tabelle ergänzen}{in 1 der letzten 5 Prüfungen}
\pfgruppe{}
\begin{pfnr}{eigene Aufgabe}{7.}{}
Die Tabelle zeigt Werte. Kreuze an, was von Wert zu Wert gleich bleibt. \wertetabelle[100,110,121]{t}{\text{Wert}}{0,1,2}
\pfkreuzzeile{die Differenz}\pfkreuzzeile{der Quotient}
\fusshilfe{\mbox{7:~der Quotient}}
\end{pfnr}
\begin{pfnr}{eigene Aufgabe}{8.}{}
Die Tabelle zeigt Werte. Kreuze an, was von Wert zu Wert gleich bleibt. \wertetabelle[100,110,120]{t}{\text{Wert}}{0,1,2}
\pfkreuzzeile{die Differenz}\pfkreuzzeile{der Quotient}
\fusshilfe{\mbox{8:~die Differenz}}
\end{pfnr}
\pfgruppe{}
\begin{pfnr}{eigene Aufgabe}{9.}{}
Ein Bestand von $300$ wächst jedes Jahr mit dem Faktor $1{,}1$. Berechne den Wert nach $4$ Jahren mit einer Potenz.
\par\noindent \leerfeld \par
\rechenplatz{2}
\fusshilfe{\mbox{9:~$439{,}23$}}
\end{pfnr}
\pfgruppe{}
\begin{pfnr}{eigene Aufgabe}{10.}{}
Eine Miete beträgt zu Beginn $300\,€$. Sie steigt jedes Jahr um $10\,\%$. Ergänze in der Tabelle das fehlende Feld. \wertetabelle[,330,363]{\text{Jahr}}{\text{Miete in €}}{0,1,2}
\rechenplatz{2}
\end{pfnr}
\begin{pfnr}{eigene Aufgabe}{11.}{}
Die Zahl der Mitglieder eines Vereins wächst jedes Jahr um $25\,\%$. Im Jahr $2$ hat der Verein $375$ Mitglieder. Wie viele Mitglieder waren es im Jahr $1$?
\par\noindent \leerfeld \par
\rechenplatz{2}
\fusshilfe{\mbox{11:~$300$}}
\end{pfnr}
\begin{pfnr}{eigene Aufgabe}{12.}{}
Die Zahl der Mitglieder ist jetzt $250$. Sie wird bei jedem Schritt mit $1{,}15$ multipliziert. Berechne den nächsten Wert.
\par\noindent \leerfeld \par
\rechenplatz{2}
\fusshilfe{\mbox{12:~$287{,}5$}}
\end{pfnr}
\begin{pfnr}{eigene Aufgabe}{13.}{}
Ein Stoff zerfällt. Nach jeder Stunde bleiben $70\,\%$ der Menge übrig. Fülle die Tabelle weiter aus. \wertetabelle[1\,000,,,]{\text{Stunde}}{\text{Menge in mg}}{0,1,2,3}
\rechenplatz{3}
\fusshilfe{\mbox{13:~$343$}}
\end{pfnr}
\begin{pfnr}{eigene Aufgabe}{14.}{}
Ein Auto kostet neu $24\,000\,€$. Es verliert jedes Jahr $15\,\%$ seines Werts. Lege eine Tabelle für die Jahre $0$ bis $3$ an. Wie viel ist das Auto nach drei Jahren noch wert?
\rechenplatz{4}
\fusshilfe{\mbox{14:~$14\,739$}}
\end{pfnr}
\begin{pfnr}{eigene Aufgabe}{15.}{}
Ein Guthaben wächst jedes Jahr um $10\,\%$. In der Zeile der Jahre fehlt eine Angabe. Zu welchem Jahr gehört der Wert $1\,331$? Prüfe mit dem Faktor. {\setlength{\mbzell}{1.7cm}\wertetabelle[1\,000,1\,100,1\,331,1\,464{,}10]{\text{Jahr}}{\text{Guthaben in €}}{0,1,?,4}}
\rechenplatz{2}
\fusshilfe{\mbox{15:~$1\,331$}}
\end{pfnr}
\pfgruppe{}
\begin{pfnr}{eigene Aufgabe}{16.}{}
Ein Bestand wächst jedes Jahr mit dem Faktor $1{,}1$. Fülle die Lücke in der Tabelle. \wertetabelle[400,440,,532{,}4]{\text{Jahr}}{\text{Anzahl}}{0,1,2,3}
\rechenplatz{3}
\fusshilfe{\mbox{16:~$484$}}
\end{pfnr}
\pfgruppe{}
\begin{pfnr}{P10 ’20}{17.}{2 BE}
Laut einer Studie aus dem Jahr 2019 verbrauchte jeder Mensch im Jahr 2019 im Schnitt 54 Tausend Liter Wasser; der Verbrauch je Person soll in Zukunft jedes Jahr um 3 \% wachsen. Ergänze die Tabelle mit diesen Angaben.
\par\smallskip\noindent \par\medskip\begingroup\centering\renewcommand{\arraystretch}{1.25}\begin{tabular}{|p{4.5cm}|p{1.6cm}|p{1.6cm}|p{1.6cm}|p{1.6cm}|}\hline Jahr & 2019 & 2020 & 2021 & 2022 \\ \hline Wasserverbrauch je Person in Tausend Litern & \leerzelle & 55,6 & 57,3 & \leerzelle \\ \hline\end{tabular}\par\endgroup\medskip\par
\rechenplatz{1}
\fusshilfe{\mbox{17:~2019: 54; 2022: 57,3 · 1,03 = 59,019 $\approx$ 59,0}}
\end{pfnr}
\pfgruppe{}
\begin{pfnr}{P10 ’16}{18.}{2 BE}
Um Krankheiten der Schilddrüse zu erkennen, spritzt man dem Patienten eine Flüssigkeit, die eine kleine Menge radioaktives Technetium enthält. Die Masse des Technetiums im Körper wird jede Stunde um etwa 11 \% kleiner. Ergänze die beiden Lücken in der Tabelle: die Masse nach 1 Stunde und die Zeit, zu der noch 3,14 mg vorhanden sind.
\par\smallskip\noindent \sachtabelle{lccccccc}{Zeit in h & 0 & 1 & 2 & \leerzelle & 5 & 6 & 8}{Masse in mg & 5,00 & \leerzelle & 3,96 & 3,14 & 2,79 & 2,48 & 1,97}\par
\rechenplatz{3}
\fusshilfe{\mbox{18:~4,45 mg; 4 h}}
\end{pfnr}
\pfgruppe{}
\begin{pfnr}{P10 ’17}{19.}{2 BE}
Auf Meereshöhe (0 km) misst man einen Luftdruck von ungefähr 1000 hPa (Hektopascal). Je höher man steigt, desto kleiner wird der Luftdruck: mit jedem Kilometer Höhe um 13 \%. Trag die beiden fehlenden Werte in die Tabelle ein.
\par\smallskip\noindent \sachtabelle{lccccccc}{Höhe in km & 0 & 1 & 2 & 3 & 5 & 8 & 10}{Luftdruck in hPa & 1000 & 870 & \leerzelle & 659 & 498 & \leerzelle & 248}\par
\rechenplatz{2}
\fusshilfe{\mbox{19:~756,9 $\approx$ 757 hPa; $\approx$ 328 hPa}}
\end{pfnr}
\pfgruppe{}
\begin{pfnr}{P10 ’19}{20.}{3 BE}
Ein Ferkel wiegt 10 kg. Aus Erfahrung weiß man, dass seine Masse jede Woche um etwa 4 \% zunimmt. Zu Beginn jeder Woche wird es gewogen. Trag die drei fehlenden Werte in die Tabelle ein.
\par\smallskip\noindent \sachtabelle{lccccc}{Zeit in Wochen & 0 & 1 & 2 & … & 5}{Masse in kg & \leerzelle & 10,400 & \leerzelle & – & \leerzelle}\par
\rechenplatz{3}
\fusshilfe{\mbox{20:~10 kg; 10,816 kg; $\approx$ 12,167 kg}}
\end{pfnr}
\pfstufekopf[spunktedarstellen]{Punkte darstellen}{in 1 der letzten 5 Prüfungen}
\pfgruppe{}
\begin{pfnr}{eigene Aufgabe}{21.}{}
Der Zahlenstrahl ist in Hunderterschritten beschriftet. Welche Zahl gehört zum Punkt $A$?
\par\smallskip\noindent \zahlenstrahl[xmin=0,xmax=600,xstep=100,karo=1]{350/A}\par
\par\noindent \leerfeld \par
\fusshilfe{\mbox{21:~$350$}}
\end{pfnr}
\begin{pfnr}{eigene Aufgabe}{22.}{}
Der Zahlenstrahl ist in Fünfzigerschritten beschriftet. Welche Zahl gehört zum Punkt $B$?
\par\smallskip\noindent \zahlenstrahl[xmin=0,xmax=400,xstep=50,karo=0.8]{225/B}\par
\par\noindent \leerfeld \par
\fusshilfe{\mbox{22:~$225$}}
\end{pfnr}
\pfgruppe{}
\begin{pfnr}{P10 ’25}{23.}{2 BE}
Man beobachtet einen Bakterienstamm, der exponentiell wächst. Die Tabelle zeigt die Zahl der Bakterien in den ersten 4 Stunden. Stelle das Wachstum der Bakterien im Koordinatensystem dar.
\par\smallskip\noindent \begin{tikzpicture}\begin{axis}[xmin=0,xmax=4,ymin=0,ymax=400,tick label style={font=\scriptsize},clip=true,/pgf/number format/use comma,axis line style={-{Stealth[length=2mm]}},every axis plot/.append style={line width=0.9pt},width=11.5cm,height=7.5cm,grid=both,grid style={mbgitter,line width=0.3pt},minor tick num=1,axis lines=left,xlabel={Zeit in Stunden},ylabel={Anzahl der Bakterien},label style={font=\small},scaled ticks=false,/pgf/number format/1000 sep={\,}]\end{axis}\end{tikzpicture}\par
\rechenplatz{1}
\end{pfnr}
\begin{pfnr}{P10 ’21}{24.}{3 BE}
Eine 40 cm hohe Kerze brennt ab. Gegeben: Tabelle Zeit in Minuten 0, 10, 20, 30, 40, 80; Höhe in cm 40, 38, 36, 34, 32, 24. Stelle diese Werte als Graph im Koordinatensystem dar. Trag dazu zuerst die fehlende Einteilung der Achsen ein.
\par\smallskip\noindent \begin{tikzpicture}[x=0.450cm,y=0.450cm]\draw[mbgitter,line width=0.3pt,step=1] (0,0) grid (26,24);\draw[line width=0.7pt,-{Stealth[length=2mm]}] (1,1) -- (25.7,1);\draw[line width=0.7pt,-{Stealth[length=2mm]}] (1,1) -- (1,23.7);\node[font=\small,anchor=north east] at (25.7,0.8) {Zeit in Minuten};\node[font=\small,anchor=south west] at (1.3,23.7) {Höhe in cm};\end{tikzpicture}\par
\rechenplatz{2}
\fusshilfe{\mbox{24:~2 cm; fallende Strecke von (0; 40) bis (80; 24)}}
\end{pfnr}
\pfgruppe{}
\begin{pfnr}{eigene Aufgabe}{25.}{}
Lies am Graphen $f$ die Werte für $t = 0$, $t = 2$ und $t = 4$ ab und trage sie in die Tabelle ein. \wertetabelle{t}{\text{Wert}}{0,2,4}
\par\smallskip\noindent \begin{ksys}[xmin=0,xmax=6,ymin=0,ymax=400,xstep=1,ystep=50,xlabel=t,ylabel=Wert,ablesen] \funktionab{50*1.4^\x}{f}{0}{6} \end{ksys}\par
\rechenplatz{4}
\fusshilfe{\mbox{25:~$192$}}
\end{pfnr}
\pfgruppe{}
\begin{pfnr}{eigene Aufgabe}{26.}{}
Die Tabelle zeigt Wertepaare. Trage jedes Wertepaar als Punkt in das Koordinatensystem ein. \wertetabelle[100,120,144,173,207]{t}{\text{Anzahl}}{0,1,2,3,4}
\par\smallskip\noindent \begin{ksys}[xmin=0,xmax=5,ymin=0,ymax=250,xstep=1,ystep=25,xlabel=t,ylabel=Anzahl]  \end{ksys}\par
\rechenplatz{1}
\end{pfnr}
\begin{pfnr}{eigene Aufgabe}{27.}{}
Im Koordinatensystem ist die senkrechte Achse in Fünfzigerschritten eingeteilt. Trage die Werte aus der Tabelle genau als Punkte ein. \wertetabelle[125,150,180,216]{t}{\text{Wert}}{0,1,2,3}
\par\smallskip\noindent \begin{ksys}[xmin=0,xmax=4,ymin=0,ymax=250,xstep=1,ystep=50,xlabel=t,ylabel=Wert]  \end{ksys}\par
\rechenplatz{2}
\end{pfnr}
\begin{pfnr}{eigene Aufgabe}{28.}{}
Im Koordinatensystem sind Punkte eingetragen. Sie zeigen ein Wachstum. Verbinde die Punkte zu einer glatten Kurve.
\par\smallskip\noindent \begin{ksys}[xmin=0,xmax=5,ymin=0,ymax=250,xstep=1,ystep=25,xlabel=t,ylabel=Anzahl] \punkt{0}{100}{} \punkt{1}{120}{} \punkt{2}{144}{} \punkt{3}{173}{} \punkt{4}{207}{} \end{ksys}\par
\rechenplatz{1}
\end{pfnr}
\begin{pfnr}{eigene Aufgabe}{29.}{}
Das Koordinatensystem hat auf jeder Achse $8$ Kästchen. Teile die Achsen so ein, dass alle Werte der Tabelle passen. Beschrifte die Achsen. Trage dann die Punkte ein. \wertetabelle[300,363,439,531]{\text{Jahr}}{\text{Euro}}{0,2,4,6}
\par\smallskip\noindent \begin{ksys}[xmin=0,xmax=8,ymin=0,ymax=8]  \end{ksys}\par
\rechenplatz{3}
\fusshilfe{\mbox{29:~Punkte eingetragen}}
\end{pfnr}
\begin{pfnr}{P10 ’16}{\llap{\textasteriskcentered\,}30.}{4 BE}
Um Krankheiten der Schilddrüse zu erkennen, spritzt man dem Patienten eine Flüssigkeit, die eine kleine Menge radioaktives Technetium enthält. Die Masse des Technetiums im Körper wird jede Stunde um etwa 11 \% kleiner. Beschrifte die Achsen des Koordinatensystems passend zur Tabelle und zeichne ein, wie die Masse abnimmt. Gegeben: Tabelle 0 h: 5,00 mg; 1 h: 4,45 mg; 2 h: 3,96 mg; 4 h: 3,14 mg; 5 h: 2,79 mg; 6 h: 2,48 mg; 8 h: 1,97 mg.
\par\smallskip\noindent \begin{tikzpicture}[x=0.450cm,y=0.450cm]\draw[mbgitter,line width=0.3pt,step=1] (0,0) grid (19,18);\draw[line width=0.7pt,-{Stealth[length=2mm]}] (1,2) -- (18.7,2);\draw[line width=0.7pt,-{Stealth[length=2mm]}] (1,2) -- (1,17.7);\draw[line width=0.5pt] (3,1.75) -- (3,2.25);\draw[line width=0.5pt] (5,1.75) -- (5,2.25);\draw[line width=0.5pt] (7,1.75) -- (7,2.25);\draw[line width=0.5pt] (9,1.75) -- (9,2.25);\draw[line width=0.5pt] (11,1.75) -- (11,2.25);\draw[line width=0.5pt] (13,1.75) -- (13,2.25);\draw[line width=0.5pt] (15,1.75) -- (15,2.25);\draw[line width=0.5pt] (17,1.75) -- (17,2.25);\draw[line width=0.5pt] (0.75,4) -- (1.25,4);\draw[line width=0.5pt] (0.75,6) -- (1.25,6);\draw[line width=0.5pt] (0.75,8) -- (1.25,8);\draw[line width=0.5pt] (0.75,10) -- (1.25,10);\draw[line width=0.5pt] (0.75,12) -- (1.25,12);\draw[line width=0.5pt] (0.75,14) -- (1.25,14);\draw[line width=0.5pt] (0.75,16) -- (1.25,16);\end{tikzpicture}\par
\rechenplatz{2}
\fusshilfe{\mbox{30:~y-Achse Masse in mg, 1 mg je Teilstrich}}
\end{pfnr}
\pfgruppe{}
\begin{pfnr}{P10 ’17}{31.}{4 BE}
Auf Meereshöhe (0 km) misst man einen Luftdruck von ungefähr 1000 hPa (Hektopascal). Je höher man steigt, desto kleiner wird der Luftdruck: mit jedem Kilometer Höhe um 13 \%. Vervollständige das Koordinatensystem und zeichne den Luftdruck in Abhängigkeit von der Höhe. Gegeben: 0 km: 1000 hPa; 1 km: 870 hPa; 2 km: 757 hPa; 3 km: 659 hPa; 5 km: 498 hPa; 8 km: 328 hPa; 10 km: 248 hPa.
\par\smallskip\noindent \begin{tikzpicture}[x=0.450cm,y=0.450cm]\draw[mbgitter,line width=0.3pt,step=1] (0,0) grid (24,24);\draw[line width=0.7pt,-{Stealth[length=2mm]}] (2,2) -- (23.7,2);\draw[line width=0.7pt,-{Stealth[length=2mm]}] (2,2) -- (2,23.7);\end{tikzpicture}\par
\rechenplatz{2}
\fusshilfe{\mbox{31:~x-Achse Höhe in km, 2 Kästchen je km}}
\end{pfnr}
\pfstufekopf[sfaktorbestimmengleic]{Faktor bestimmen, Gleichung aufstellen}{in 2 der letzten 5 Prüfungen}
\pfgruppe{}
\begin{pfnr}{eigene Aufgabe}{32.}{}
Ein Wert wird jedes Jahr mit $1{,}35$ multipliziert. Kreuze an, ob der Wert mehr oder weniger wird. 
\pfkreuzzeile{mehr}\pfkreuzzeile{weniger}
\fusshilfe{\mbox{32:~der Faktor ist größer als eins}}
\end{pfnr}
\begin{pfnr}{eigene Aufgabe}{33.}{}
Ein Bestand wird ab $2016$ jedes Jahr größer. Kreuze an, welcher Schritt $t$ zum Jahr $2021$ gehört. 
\pfkreuzzeile{$t = 2021$}\pfkreuzzeile{$t = 5$}\pfkreuzzeile{$t = 6$}
\fusshilfe{\mbox{33:~$2016$ ist der Schritt null}}
\end{pfnr}
\pfgruppe{}
\begin{pfnr}{eigene Aufgabe}{34.}{}
Ein Wert steigt jedes Jahr um $40\,\%$. Mit welchem Faktor wird er jedes Jahr multipliziert?
\par\noindent Faktor \leerfeld \par
\fusshilfe{\mbox{34:~$1{,}4$}}
\end{pfnr}
\pfgruppe{}
\begin{pfnr}{eigene Aufgabe}{35.}{}
Ein Roller ist $3\,000\,€$ wert. Er verliert jedes Jahr $20\,\%$ seines Werts. Berechne Schritt für Schritt den Wert nach $1$, $2$ und $3$ Jahren.
\par\noindent \leerfeld , \leerfeld , \leerfeld \par
\rechenplatz{3}
\fusshilfe{\mbox{35:~$1\,536$}}
\end{pfnr}
\begin{pfnr}{eigene Aufgabe}{36.}{}
Im Text steht: „Auf einem Blatt sitzen $40$ Blattläuse. Ihre Zahl wächst jeden Tag um $30\,\%$.“ Unterstreiche im Text den Anfangswert und den Prozentsatz. Gib dann den Anfangswert und den Faktor an.
\par\noindent Anfangswert \leerfeld , Faktor \leerfeld \par
\rechenplatz{2}
\fusshilfe{\mbox{36:~$40$ und $1{,}3$}}
\end{pfnr}
\pfgruppe{}
\begin{pfnr}{eigene Aufgabe}{37.}{}
Ein Bestand von $2\,000$ Tieren nimmt jedes Jahr um $10\,\%$ ab. Kreuze an, welche Gleichung den Bestand $B$ nach $t$ Jahren beschreibt. 
\pfkreuzzeile{$B(t) = 2\,000 \cdot 1{,}1^t$}\pfkreuzzeile{$B(t) = 2\,000 \cdot 0{,}9^t$}\pfkreuzzeile{$B(t) = 2\,000 - 10t$}\pfkreuzzeile{$B(t) = 2\,000 \cdot 0{,}1^t$}
\fusshilfe{\mbox{37:~Abnahme um $10\,\%$ heißt Faktor $0{,}9$}}
\end{pfnr}
\pfgruppe{}
\begin{pfnr}{eigene Aufgabe}{38.}{}
In einem Becher sind $600\,\mathrm{mg}$ eines Stoffs. Jede Stunde zerfallen $20\,\%$ davon. Stelle die Gleichung für die Menge $m$ nach $t$ Stunden auf.
\par\noindent m(t) = \leerfeld \par
\rechenplatz{2}
\fusshilfe{\mbox{38:~$600 \cdot 0{,}8^t$}}
\end{pfnr}
\begin{pfnr}{eigene Aufgabe}{39.}{}
Ein Roller ist $3\,000\,€$ wert. Er verliert jedes Jahr $20\,\%$ seines Werts. Berechne den Wert nach $3$ Jahren mit einer Potenz.
\par\noindent \leerfeld[€] \par
\rechenplatz{2}
\fusshilfe{\mbox{39:~$1\,536\,€$}}
\end{pfnr}
\begin{pfnr}{eigene Aufgabe}{40.}{}
Die Menge eines Stoffs sinkt in zwei Stunden von $1\,000\,\mathrm{mg}$ auf $810\,\mathrm{mg}$. Um wie viel Prozent nimmt die Menge in jeder Stunde ab?
\par\noindent um \leerfeld[\%] \par
\rechenplatz{2}
\fusshilfe{\mbox{40:~$0{,}9^2$, Faktor $0{,}9$}}
\end{pfnr}
\begin{pfnr}{eigene Aufgabe}{41.}{}
Auf einem Blatt sitzen $40$ Blattläuse. Ihre Zahl wächst jeden Tag um $30\,\%$. Stelle die Gleichung für die Anzahl $A$ nach $t$ Tagen auf.
\par\noindent A(t) = \leerfeld \par
\rechenplatz{2}
\fusshilfe{\mbox{41:~$40 \cdot 1{,}3^t$}}
\end{pfnr}
\begin{pfnr}{eigene Aufgabe}{42.}{}
In einer Tabelle stehen nacheinander die Werte $350$ und $420$. Berechne den Faktor. Gib auch den Prozentsatz an.
\par\noindent Faktor \leerfeld , \leerfeld[\%] \par
\rechenplatz{2}
\fusshilfe{\mbox{42:~$1{,}2$, also $20\,\%$ Zunahme}}
\end{pfnr}
\pfgruppe{}
\begin{pfnr}{eigene Aufgabe}{43.}{}
Ein Wert wird jedes Jahr mit $1{,}15$ multipliziert. Um wie viel Prozent nimmt er jährlich zu?
\par\noindent um \leerfeld[\%] \par
\fusshilfe{\mbox{43:~Zunahme um $15\,\%$}}
\end{pfnr}
\pfgruppe{}
\begin{pfnr}{eigene Aufgabe}{44.}{}
Die Tabelle zeigt Werte. Kreuze an, zu welcher Gleichung die Tabelle gehört. \wertetabelle[80,100,125]{t}{\text{Wert}}{0,1,2}
\pfkreuzzeile{$f(t) = 80 \cdot 1{,}25^t$}\pfkreuzzeile{$f(t) = 80 + 20t$}\pfkreuzzeile{$f(t) = 80 \cdot 0{,}8^t$}
\fusshilfe{\mbox{44:~Startwert $80$ und Quotient $1{,}25$}}
\end{pfnr}
\pfgruppe{}
\begin{pfnr}{eigene Aufgabe}{45.}{}
Die Menge eines Medikaments nimmt jede Stunde um $25\,\%$ ab. Mit welchem Faktor rechnest du die Menge in der nächsten Stunde aus?
\par\noindent Faktor \leerfeld \par
\fusshilfe{\mbox{45:~$0{,}75$}}
\end{pfnr}
\begin{pfnr}{eigene Aufgabe}{46.}{}
Die Zahl der Mitglieder eines Vereins wird mit $M(t) = 250 \cdot 1{,}05^t$ berechnet. Erkläre, was $250$, $1{,}05$ und $t$ bedeuten.
\rechenplatz{2}
\fusshilfe{\mbox{46:~$1{,}05$: Faktor, jedes Jahr $5\,\%$ mehr}}
\end{pfnr}
\pfgruppe{}
\begin{pfnr}{P10 ’17}{\llap{\textasteriskcentered\,}47.}{1 BE}
Auf Meereshöhe (0 km) misst man einen Luftdruck von ungefähr 1000 hPa (Hektopascal). Je höher man steigt, desto kleiner wird der Luftdruck: mit jedem Kilometer Höhe um 13 \%. Die Abnahme des Luftdrucks lässt sich mit einer Gleichung beschreiben ($x$: Höhe in km, $y$: Luftdruck in hPa). Kreuze die richtige an:
\pfkreuzzeile{$y = 0{,}87 \cdot x$}\pfkreuzzeile{$y = 1000 \cdot 0{,}87^x$}\pfkreuzzeile{$y = x^{1{,}13}$}\pfkreuzzeile{$y = 1000 \cdot 1{,}13^x$}
\par\smallskip\noindent \sachtabelle{lccccccc}{Höhe in km & 0 & 1 & 2 & 3 & 5 & 8 & 10}{Luftdruck in hPa & 1000 & 870 & \leerzelle & 659 & 498 & \leerzelle & 248}\par
\fusshilfe{\mbox{47:~y = 1000 · $0{,}87^{x}$}}
\end{pfnr}
\pfgruppe{}
\begin{pfnr}{P10 ’18}{48.}{2 BE}
Ein junges Berliner Unternehmen zeigt in einer Tabelle, wie sein Umsatz gewachsen ist: 2013: 600 000 €, 2014: 648 000 €, 2015: 699 840 €, 2016: 755 827 €, 2017: 816 293 €. Es behauptet, der Umsatz wachse jedes Jahr um 8 \%. Zeige, dass das für 2014 stimmt.
\par\smallskip\noindent \sachtabelle{lrrrrr}{Jahr & 2013 & 2014 & 2015 & 2016 & 2017}{Umsatz in € & 600 000 & 648 000 & 699 840 & 755 827 & 816 293}\par
\rechenplatz{1}
\fusshilfe{\mbox{48:~600 000 € · 1,08 = 648 000 €}}
\end{pfnr}
\begin{pfnr}{P10 ’19}{\llap{\textasteriskcentered\,}49.}{3 BE}
Ein Ferkel wiegt 10 kg. Aus Erfahrung weiß man, dass seine Masse jede Woche um etwa 4 \% zunimmt. Zu Beginn jeder Woche wird es gewogen. Die Bäuerin Wilma berechnet die Masse des Ferkels nach $x$ Wochen mit $f(x) = 10 \cdot 1{,}04^x$. Schreib auf, was $10$, $1{,}04$ und $x$ in dieser Gleichung bedeuten.
\par\smallskip\noindent \sachtabelle{lc}{Teil der Gleichung & Bedeutung}{10 & \leerzelle\\ 1,04 & \leerzelle\\ x & \leerzelle}\par
\rechenplatz{1}
\fusshilfe{\mbox{49:~1,04: Wachstumsfaktor, +4 \% je Woche}}
\end{pfnr}
\begin{pfnr}{eigene Aufgabe}{50.}{}
Die Tabelle zeigt, wie sich eine Anzahl entwickelt. Berechne die Quotienten von Wert zu Wert. Teile dazu immer den neuen Wert durch den alten. Schreibe das Ergebnis in einem Satz auf. \wertetabelle[150,180,216,259{,}2]{\text{Jahr}}{\text{Anzahl}}{0,1,2,3}
\par\noindent Quotienten: \leerfeld , \leerfeld , \leerfeld \par
\rechenplatz{4}
\end{pfnr}
\pfgruppe{}
\begin{pfnr}{eigene Aufgabe}{51.}{}
Eine Gemeinde hatte im Jahr $2019$ genau $2\,500$ Einwohner. Die Zahl wächst jedes Jahr um $6\,\%$. Wie viele Einwohner sind es im Jahr $2025$? Runde auf ganze Personen.
\par\noindent t = \leerfeld , $\approx$ \leerfeld \par
\rechenplatz{2}
\fusshilfe{\mbox{51:~$\approx 3\,546$}}
\end{pfnr}
\begin{pfnr}{eigene Aufgabe}{52.}{}
Ein Guthaben von $2\,000\,€$ wächst jedes Jahr um $4\,\%$. Jemand sagt: „Nach zehn Jahren sind es über $3\,000\,€$.“ Prüfe mit einer Rechnung, ob das stimmt.
\rechenplatz{2}
\end{pfnr}
\pfgruppe{}
\begin{pfnr}{P10 ’17}{53.}{2 BE}
Auf Meereshöhe (0 km) misst man einen Luftdruck von ungefähr 1000 hPa (Hektopascal). Je höher man steigt, desto kleiner wird der Luftdruck: mit jedem Kilometer Höhe um 13 \%. Eine Bergsteigerin misst 573 hPa. Entscheide, ob sie ungefähr 4 km hoch sein kann, und begründe.
\par\smallskip\noindent \sachtabelle{lccccccc}{Höhe in km & 0 & 1 & 2 & 3 & 5 & 8 & 10}{Luftdruck in hPa & 1000 & 870 & \leerzelle & 659 & 498 & \leerzelle & 248}\par
\rechenplatz{1}
\fusshilfe{\mbox{53:~ja, 1000 · 0,87$^{4}$ $\approx$ 573 hPa}}
\end{pfnr}
\pfgruppe{}
\begin{pfnr}{P10 ’16}{\llap{\textasteriskcentered\,}54.}{1 BE}
Um Krankheiten der Schilddrüse zu erkennen, spritzt man dem Patienten eine Flüssigkeit, die eine kleine Menge radioaktives Technetium enthält. Die Masse des Technetiums im Körper wird jede Stunde um etwa 11 \% kleiner. Berechne, wie viele Milligramm Technetium 24 Stunden nach der Spritze noch im Körper sind.
\par\smallskip\noindent \sachtabelle{lccccccc}{Zeit in h & 0 & 1 & 2 & \leerzelle & 5 & 6 & 8}{Masse in mg & 5,00 & \leerzelle & 3,96 & 3,14 & 2,79 & 2,48 & 1,97}\par
\rechenplatz{2}
\fusshilfe{\mbox{54:~$\approx$ 0,31 mg}}
\end{pfnr}
\pfgruppe{}
\begin{pfnr}{P10 ’20}{\llap{\textasteriskcentered\,}55.}{3 BE}
Nach einer Studie lag der Wasserverbrauch 2019 bei durchschnittlich 54 Tausend Litern pro Person und Jahr; er soll jedes Jahr um 3 \% steigen. Diese Vorhersage beschreibt die Gleichung $y = 54\cdot 1{,}03^{x}$ ($x$: Jahre ab 2019, $y$: Verbrauch pro Person in Tausend Litern). Eine zweite Studie erwartet, dass der Verbrauch pro Person von 2019 bis 2033 insgesamt um etwa 40 \% zunimmt. Untersuche, ob beide Studien für 2033 denselben Verbrauch vorhersagen.
\rechenplatz{3}
\fusshilfe{\mbox{55:~75,6}}
\end{pfnr}
\pfgruppe{}
\begin{pfnr}{P10 ’25}{\llap{\textasteriskcentered\,}56.}{4 BE}
Man beobachtet einen Bakterienstamm, der exponentiell wächst. Die Tabelle zeigt die Zahl der Bakterien in den ersten 4 Stunden. Zeige, dass der Wachstumsfaktor 1,4 ist. Gib eine Gleichung an, die die Zahl der Bakterien nach $t$ Stunden beschreibt. Prüfe mit einer Rechnung, ob diese Aussage stimmt: „Nach 24 Stunden sind es mehr als 250 000 Bakterien.“
\par\smallskip\noindent \sachtabelle{lrrrrr}{Zeit in Stunden & 0 & 1 & 2 & 3 & 4}{Anzahl der Bakterien & 80 & 112 & 157 & 220 & 308}\par
\rechenplatz{4}
\fusshilfe{\mbox{56:~1,4; 80 · $1{,}4^{t}$; $\approx$ 257 136 > 250 000}}
\end{pfnr}
\pfstufekopf[sgraphzuordnenundbegr]{Graph zuordnen und begründen}{in 1 der letzten 5 Prüfungen}
\pfgruppe{}
\begin{pfnr}{eigene Aufgabe}{57.}{}
Sieh dir im Koordinatensystem den Graphen $f$ an. Kreuze an, bei welchem Wert er die senkrechte Achse schneidet. 
\pfkreuzzeile{$0$}\pfkreuzzeile{$50$}\pfkreuzzeile{$100$}
\par\smallskip\noindent \begin{ksys}[xmin=0,xmax=10,ymin=0,ymax=700,xstep=1,ystep=100,xlabel=t in Jahren,ylabel=Wert,ablesen] \funktionab{100*1.2^\x}{f}{0}{10} \gerade{30}{100}{g} \parabel{6}{0}{0}{h} \end{ksys}\par
\fusshilfe{\mbox{57:~$100$}}
\end{pfnr}
\begin{pfnr}{eigene Aufgabe}{58.}{}
Das Koordinatensystem zeigt drei Graphen $f$, $g$ und $h$. Kreuze an, welcher Graph im Ursprung beginnt, also ganz unten links. 
\pfkreuzzeile{$f$}\pfkreuzzeile{$g$}\pfkreuzzeile{$h$}
\par\smallskip\noindent \begin{ksys}[xmin=0,xmax=10,ymin=0,ymax=1000,xstep=1,ystep=100,xlabel=t in Jahren,ylabel=Wert,ablesen] \funktionab{240*1.15^\x}{f}{0}{10} \gerade{48}{240}{g} \parabel{8}{0}{0}{h} \end{ksys}\par
\fusshilfe{\mbox{58:~nur er geht durch $(0 | 0)$}}
\end{pfnr}
\pfgruppe{}
\begin{pfnr}{eigene Aufgabe}{59.}{}
Eine Tasse Tee kühlt ab. Der Tee ist wärmer als die Luft im Raum. Dieser Unterschied wird jede Minute um $10\,\%$ kleiner. Kreuze an, welcher Graph im Koordinatensystem zu dem Unterschied passt. 
\pfkreuzzeile{$f$}\pfkreuzzeile{$g$}\pfkreuzzeile{$h$}
\par\smallskip\noindent \begin{ksys}[xmin=0,xmax=10,ymin=0,ymax=100,xstep=1,ystep=10,xlabel=t in Jahren,ylabel=Wert,ablesen] \funktionab{60*0.9^\x}{f}{0}{10} \gerade{-5}{60}{g} \funktionab{60*1.1^\x}{h}{0}{10} \end{ksys}\par
\fusshilfe{\mbox{59:~fallende Kurve, die immer flacher wird}}
\end{pfnr}
\pfgruppe{}
\begin{pfnr}{eigene Aufgabe}{60.}{}
Ein Pilz bedeckt am Anfang eine Fläche von $100\,\mathrm{cm}^2$. Die Fläche wächst jeden Tag um $20\,\%$. Das Koordinatensystem zeigt drei Graphen $f$, $g$ und $h$. Begründe für jeden Graphen in einem Satz, ob er passt.
\par\smallskip\noindent \begin{ksys}[xmin=0,xmax=10,ymin=0,ymax=700,xstep=1,ystep=100,xlabel=t in Tagen,ylabel=Fläche,ablesen] \funktionab{100*1.2^\x}{f}{0}{10} \gerade{30}{100}{g} \parabel{6}{0}{0}{h} \end{ksys}\par
\rechenplatz{1}
\end{pfnr}
\pfgruppe{}
\begin{pfnr}{eigene Aufgabe}{61.}{}
In einem Wasserbecken sind $300$ Liter. Jede Minute fließen $60$ Liter dazu. Kreuze an, welcher Graph im Koordinatensystem zur Wassermenge passt. 
\pfkreuzzeile{$f$}\pfkreuzzeile{$g$}\pfkreuzzeile{$h$}
\par\smallskip\noindent \begin{ksys}[xmin=0,xmax=10,ymin=0,ymax=1000,xstep=1,ystep=100,xlabel=t in min,ylabel=Liter,ablesen] \funktionab{300*1.12^\x}{f}{0}{10} \gerade{60}{300}{g} \parabel{9}{0}{0}{h} \end{ksys}\par
\end{pfnr}
\pfgruppe{}
\begin{pfnr}{eigene Aufgabe}{\llap{\textasteriskcentered\,}62.}{3 BE}
Ein Kalb wiegt $45\,\mathrm{kg}$. Es nimmt jede Woche um $5\,\%$ zu. Das Koordinatensystem zeigt drei Graphen $f$, $g$ und $h$. Entscheide für jeden Graphen, ob er passt. Begründe jede Entscheidung.
\par\smallskip\noindent \begin{ksys}[xmin=0,xmax=10,ymin=0,ymax=80,xstep=1,ystep=10,xlabel=t in Wochen,ylabel=kg,ablesen] \funktionab{45*1.05^\x}{f}{0}{10} \gerade{2}{45}{g} \parabel{0.7}{0}{0}{h} \end{ksys}\par
\rechenplatz{1}
\end{pfnr}
\begin{pfnr}{eigene Aufgabe}{63.}{}
Das Koordinatensystem zeigt drei Graphen $a$, $b$ und $c$. Welcher Graph ist eine Gerade?
\par\smallskip\noindent \begin{ksys}[xmin=0,xmax=10,ymin=0,ymax=1000,xstep=1,ystep=100,xlabel=t in Jahren,ylabel=Wert,ablesen] \funktionab{250*1.14^\x}{a}{0}{10} \gerade{45}{250}{b} \parabel{8}{0}{0}{c} \end{ksys}\par
\par\noindent \leerfeld \par
\fusshilfe{\mbox{63:~$b$ ist die Gerade}}
\end{pfnr}
\begin{pfnr}{eigene Aufgabe}{64.}{}
Ein Konto wächst jedes Jahr um $15\,\%$. Das Koordinatensystem zeigt drei Graphen. Begründe in einem Satz, warum der Graph $g$ nicht zu dem Konto passt.
\par\smallskip\noindent \begin{ksys}[xmin=0,xmax=10,ymin=0,ymax=1000,xstep=1,ystep=100,xlabel=t in Jahren,ylabel=Wert,ablesen] \funktionab{240*1.15^\x}{f}{0}{10} \gerade{48}{240}{g} \parabel{8}{0}{0}{h} \end{ksys}\par
\rechenplatz{2}
\fusshilfe{\mbox{64:~eine Gerade hat immer denselben Zuwachs}}
\end{pfnr}
\begin{pfnr}{P10 ’19}{\llap{\textasteriskcentered\,}65.}{3 BE}
Ein Ferkel wiegt 10 kg. Aus Erfahrung weiß man, dass seine Masse jede Woche um etwa 4 \% zunimmt. Zu Beginn jeder Woche wird es gewogen. Für jeden der drei Graphen: Entscheide, ob er zum Wachstum des Ferkels passt, und begründe.
\par\smallskip\noindent \begin{tabular}[t]{@{}c@{}}{\small 1}\enspace$\square$\\\begin{tikzpicture}\begin{axis}[width=3.9cm,height=3.4cm,scale only axis=false,axis lines=middle,xtick=\empty,ytick=\empty,xmin=0,xmax=1.08,ymin=0,ymax=1.08,clip=true,axis line style={-{Stealth[length=1.5mm]}},grid=both,minor tick num=0,xtick={0.2,0.4,0.6,0.8,1},ytick={0.2,0.4,0.6,0.8,1},xticklabels={},yticklabels={},grid style={mbgitter},xlabel={\tiny Zeit in Wochen},ylabel={\tiny Masse in kg},x label style={anchor=north east,at={(1,0)},yshift=-2pt},y label style={anchor=south west,at={(0,1)},xshift=1pt}]\addplot[black,line width=0.9pt,domain=0:1,samples=80] {0.25+0.6*x};\end{axis}\end{tikzpicture}\\[1pt]{\footnotesize $\square$ passt\quad $\square$ passt nicht}\\[2pt]\rule{3cm}{0.4pt}\end{tabular}\hspace{0.4cm}\begin{tabular}[t]{@{}c@{}}{\small 2}\enspace$\square$\\\begin{tikzpicture}\begin{axis}[width=3.9cm,height=3.4cm,scale only axis=false,axis lines=middle,xtick=\empty,ytick=\empty,xmin=0,xmax=1.08,ymin=0,ymax=1.08,clip=true,axis line style={-{Stealth[length=1.5mm]}},grid=both,minor tick num=0,xtick={0.2,0.4,0.6,0.8,1},ytick={0.2,0.4,0.6,0.8,1},xticklabels={},yticklabels={},grid style={mbgitter},xlabel={\tiny Zeit in Wochen},ylabel={\tiny Masse in kg},x label style={anchor=north east,at={(1,0)},yshift=-2pt},y label style={anchor=south west,at={(0,1)},xshift=1pt}]\addplot[black,line width=0.9pt,domain=0:1,samples=80] {0.2*exp(1.45*x)};\end{axis}\end{tikzpicture}\\[1pt]{\footnotesize $\square$ passt\quad $\square$ passt nicht}\\[2pt]\rule{3cm}{0.4pt}\end{tabular}\hspace{0.4cm}\begin{tabular}[t]{@{}c@{}}{\small 3}\enspace$\square$\\\begin{tikzpicture}\begin{axis}[width=3.9cm,height=3.4cm,scale only axis=false,axis lines=middle,xtick=\empty,ytick=\empty,xmin=0,xmax=1.08,ymin=0,ymax=1.08,clip=true,axis line style={-{Stealth[length=1.5mm]}},grid=both,minor tick num=0,xtick={0.2,0.4,0.6,0.8,1},ytick={0.2,0.4,0.6,0.8,1},xticklabels={},yticklabels={},grid style={mbgitter},xlabel={\tiny Zeit in Wochen},ylabel={\tiny Masse in kg},x label style={anchor=north east,at={(1,0)},yshift=-2pt},y label style={anchor=south west,at={(0,1)},xshift=1pt}]\addplot[black,line width=0.9pt,domain=0:1,samples=80] {0.88*(x^2)};\end{axis}\end{tikzpicture}\\[1pt]{\footnotesize $\square$ passt\quad $\square$ passt nicht}\\[2pt]\rule{3cm}{0.4pt}\end{tabular}\par
\rechenplatz{1}
\fusshilfe{\mbox{65:~Graph 1 nein; Graph 2 ja; Graph 3 nein}}
\end{pfnr}
\pfpruefstein{P10 ’26}
\begin{pfnr}{}{}{}
Can bezieht Anfang 2026 eine neue Wohnung. Laut Mietvertrag steigt die Miete jedes Jahr um 1,9 \%. Im ersten Jahr zahlt er 650 € im Monat.
\end{pfnr}
\begin{pfnr}{}{a)}{2 BE}
Can legt eine Tabelle an. Vervollständige sie.
\par\smallskip\noindent \sachtabelle{lcccc}{Jahr & 2026 & 2027 & 2028 & 2029}{Miete in € & \leerzelle & 662,35 & 674,93 & \leerzelle}\par
\rechenplatz{3}
\end{pfnr}
\begin{pfnr}{}{b)}{3 BE}
Entscheide, welcher Graph zur Entwicklung der Miete passt, und kreuze ihn an. Erkläre bei den zwei übrigen Graphen jeweils, weshalb sie nicht passen.
\par\smallskip\noindent \begin{tabular}[t]{@{}c@{}}{\small A}\enspace$\square$\\\begin{tikzpicture}\begin{axis}[width=3.9cm,height=3.4cm,scale only axis=false,axis lines=middle,xtick=\empty,ytick=\empty,xmin=0,xmax=1.08,ymin=0,ymax=1.08,clip=true,axis line style={-{Stealth[length=1.5mm]}},grid=both,minor tick num=0,xtick={0.2,0.4,0.6,0.8,1},ytick={0.2,0.4,0.6,0.8,1},xticklabels={},yticklabels={},grid style={mbgitter},xlabel={\tiny Zeit in Jahren},ylabel={\tiny Miete in €},x label style={anchor=north east,at={(1,0)},yshift=-2pt},y label style={anchor=south west,at={(0,1)},xshift=1pt}]\addplot[black,line width=0.9pt,domain=0:1,samples=80] {0.88*(x^2)};\end{axis}\end{tikzpicture}\end{tabular}\hspace{0.4cm}\begin{tabular}[t]{@{}c@{}}{\small B}\enspace$\square$\\\begin{tikzpicture}\begin{axis}[width=3.9cm,height=3.4cm,scale only axis=false,axis lines=middle,xtick=\empty,ytick=\empty,xmin=0,xmax=1.08,ymin=0,ymax=1.08,clip=true,axis line style={-{Stealth[length=1.5mm]}},grid=both,minor tick num=0,xtick={0.2,0.4,0.6,0.8,1},ytick={0.2,0.4,0.6,0.8,1},xticklabels={},yticklabels={},grid style={mbgitter},xlabel={\tiny Zeit in Jahren},ylabel={\tiny Miete in €},x label style={anchor=north east,at={(1,0)},yshift=-2pt},y label style={anchor=south west,at={(0,1)},xshift=1pt}]\addplot[black,line width=0.9pt,domain=0:1,samples=80] {0.2*exp(1.45*x)};\end{axis}\end{tikzpicture}\end{tabular}\hspace{0.4cm}\begin{tabular}[t]{@{}c@{}}{\small C}\enspace$\square$\\\begin{tikzpicture}\begin{axis}[width=3.9cm,height=3.4cm,scale only axis=false,axis lines=middle,xtick=\empty,ytick=\empty,xmin=0,xmax=1.08,ymin=0,ymax=1.08,clip=true,axis line style={-{Stealth[length=1.5mm]}},grid=both,minor tick num=0,xtick={0.2,0.4,0.6,0.8,1},ytick={0.2,0.4,0.6,0.8,1},xticklabels={},yticklabels={},grid style={mbgitter},xlabel={\tiny Zeit in Jahren},ylabel={\tiny Miete in €},x label style={anchor=north east,at={(1,0)},yshift=-2pt},y label style={anchor=south west,at={(0,1)},xshift=1pt}]\addplot[black,line width=0.9pt,domain=0:1,samples=80] {0.25+0.6*x};\end{axis}\end{tikzpicture}\end{tabular}\par
\rechenplatz{3}
\end{pfnr}
\begin{pfnr}{}{*c)}{3 BE}
Stelle eine Gleichung auf, mit der man die Monatsmiete nach beliebig vielen Jahren berechnen kann. Bestimme die Monatsmiete im Jahr 2040.
\rechenplatz{3}
\end{pfnr}
\end{document}